\section{State of the Art}

Recent advances in the literature on algorithmic fairness have refined and expanded the set of available strategies for bias mitigation in machine learning. Mitigation methods are typically classified into three categories: pre-processing, in-processing, and post-processing. Pre-processing methods focus on adjusting the dataset prior to training, aiming to neutralize unwanted correlations between sensitive attributes and the target variable. One prominent example is counterfactual data augmentation, where synthetic instances are generated with swapped sensitive attributes to reduce statistical dependence between protected variables and model outputs~\cite{pessach2020}. In-processing techniques incorporate fairness constraints directly into the learning process. A notable approach is representation neutralization, which modifies the internal representations learned by the model to obscure sensitive attributes, enhancing fairness without significant losses in predictive performance~\cite{caton2020}. Post-processing methods operate after training, altering the final predictions to meet fairness criteria. These techniques are especially useful when access to training data is restricted or when model retraining is not feasible. Recent proposals include calibrated adjustment algorithms that re-score outputs to ensure fairness metrics such as demographic parity and equal opportunity are satisfied~\cite{mehrabi2021}.

In the financial domain, particularly in credit scoring, the deployment of fairness-aware machine learning models is both a technical and regulatory challenge. Disparities in credit approval rates across demographic groups—such as age, gender, or nationality—may stem from historical bias encoded in training data. These disparities, when left unmitigated, not only compromise fairness but may also violate fair lending regulations~\cite{fuster2020}. Consequently, recent studies have evaluated fairness mitigation methods specifically in credit contexts, balancing predictive performance with ethical and legal considerations. For example, benchmark experiments with reweighing and reject option classification have demonstrated consistent reductions in demographic disparities without degrading model accuracy~\cite{bellamy2019}. Furthermore, the literature has acknowledged the role of proxy variables, which are non-sensitive features that correlate with protected attributes and may inadvertently reintroduce bias even when direct sensitive attributes are excluded. As the use of AI models in financial decision-making expands, ensuring transparency and fairness through systematic auditing has become a central concern, especially for institutions operating under strict regulatory frameworks.
